\section{The Proposed Solution: CarbonPass}\label{sec:solution}

\textbf{CarbonPass is a local-first AI compliance and energy-optimization
copilot for export SMEs, prototyped on the fastener cluster.} In one sentence
of user experience: \emph{photograph your electricity bill and invoices inside
LINE, answer three questions in Mandarin or Taiwanese, and receive (a) the
carbon data pack your European buyer needs, (b) the money it saves them, and
(c) a production schedule that cuts your power bill and makes next quarter's
numbers better.}

\begin{table}[htbp]
\caption{The CarbonPass pipeline at a glance}
\label{tab:pipeline}
\small
\begin{tabularx}{\textwidth}{@{}p{0.2\textwidth}X@{}}
\toprule
\textbf{Stage} & \textbf{What happens} \\
\midrule
Inputs the SME already has & Taipower bills, material invoices (steel wire rod, fuels), production records --- photographed in LINE, in Mandarin/Taiwanese \\
Module 1: Data-Pack Engine & Document AI + allocation across product lines $\to$ per-CN-code embedded emissions in the EU template, flagged field-by-field for the verifier \\
Module 2: Grid-Aware Scheduler & Hourly grid carbon intensity from open 10-minute Taipower data + time-of-use tariffs $\to$ shift plan for furnaces and compressors; NT\$ and tCO$_2$e deltas \\
Outputs & CBAM data pack; organizational inventory; ``default vs.\ actual'' buyer-cost screen; monitored reduction ledger \\
\bottomrule
\end{tabularx}
\end{table}

\subsection{Module 1 --- the Data-Pack Engine (compliance core)}

Document intelligence ingests bills, invoices and production records
\textbf{on the factory's own device or premises} --- raw data never leaves. The
engine then performs the two computations generic calculators cannot:
\textbf{allocation} (splitting facility energy across product lines and CN
codes using production volumes, machine ratings and operating-hour heuristics,
with quantified uncertainty per line --- a strict CBAM requirement that
invoices alone cannot satisfy) and \textbf{assembly of ``specific embedded
emissions'' per the EU rulebook} (simple-vs-complex-goods logic, wire-rod
precursor from mill certificates or conservative default fallback, output
rendered in the European Commission's template structure). The most persuasive
output is one screen: \emph{``with your data, your buyer surrenders
$\approx$EUR~150/t; without it, $\approx$EUR~450--750/t and rising''} --- a
retention argument the owner can forward to the customer.

\subsection{Module 2 --- the Grid-Aware Scheduler (environmental-economic engine)}

Taiwan publishes its power system's \textbf{generation by unit every 10 minutes
as open data}\cite{taipowerdata}. CarbonPass computes an hourly
grid-carbon-intensity curve from that public feed, overlays Taipower
time-of-use tariffs (summer peak rates run several times the off-peak
rate\cite{taipower}), and runs a scheduling optimizer over the factory's
\emph{flexible} loads --- heat-treatment furnaces, compressors, plating
rectifiers --- against order deadlines stated in the chat. Input is 15-minute
smart-meter data the customer exports from their own Taipower account
(consent-based; no third-party API dependency). Output: a shift plan, projected
NT\$ savings, projected tCO$_2$e reduction, and a monitored before/after ledger
--- a continuously documented reduction trajectory structured to the
baseline-and-monitoring logic of ISO~14064-2. Stated carefully:
\emph{structured to the standard's logic; certification remains the accredited
verifier's act.} Every avoided tonne also shrinks the next data pack ---
measurement and mitigation feed each other.

\subsection{Module 3 --- the Inclusion Interface (the reason this belongs in this hackathon)}

The front end is not a dashboard but a \textbf{conversation in LINE,
Mandarin/Taiwanese-first, voice-friendly, zero ESG vocabulary}.
``{\cjk 拍一張電費單}'' --- \emph{snap a photo of the power bill} --- is the
entire onboarding. The excluded party here is a person: a 60-plus
owner-operator who will never log into a European portal or read a 2,400-page
implementing regulation\cite{omm}, and should not have to. Digital inclusion
--- the Handbook's ``ages, regions, and needs'' --- is delivered as an
interface property, not a slogan.

\subsection{One engine, both directions of readiness}

\begin{table}[htbp]
\caption{Extended uses of the same data spine}
\label{tab:extended}
\small
\begin{tabularx}{\textwidth}{@{}p{0.17\textwidth}Xp{0.3\textwidth}@{}}
\toprule
\textbf{Direction} & \textbf{Artifact / use} & \textbf{Trigger} \\
\midrule
International & CBAM data packs per CN code & Definitive phase, live since Jan 2026\cite{taxud} \\
 & Customer Scope-3 questionnaires (TSMC-chain, global brands) & SBTi/RE100 supplier programs\cite{tsmc} \\
 & ISO 14067 preparation; Digital Product Passport payloads & ESPR phase-in 2027--2030\cite{espr} \\
\addlinespace
Domestic & Standing organizational inventory (ISO 14064-1-style) & Carbon-fee threshold lowering\cite{ecobiz} \\
 & Documented reduction ledger, positioned for the voluntary-credit mechanism (credits offset fee liability at 1.2:1, capped at 10\%) & MOENV voluntary reduction rules\cite{offset} \\
 & Bank-ready carbon data for sustainability-linked lending & FSC Green Finance policies \\
\bottomrule
\end{tabularx}
\end{table}

\section{Beyond Today's Scope: The Mandates That Extend the Idea}\label{sec:future}

A proposal must justify its life after the demo. CarbonPass's forward case
rests on dated regulatory trajectories, not speculation:

\begin{enumerate}[leftmargin=1.4em]
\item \textbf{The CBAM downstream extension (decision imminent).} $\sim$180
added product categories, obligations from 1~January~2028; Council position
adopted June 2026; \textbf{European Parliament plenary vote expected September
2026 --- inside this Hackathon's mentorship window}\cite{akin}. By the final
review in late October, the team presents with the outcome known: either the
addressable market grows from 2,600 firms to the exporters of 180 categories,
or the beachhead is confirmed with more preparation time. Treated correctly as
pending legislation, this is a roadmap event few proposals anywhere can
schedule between submission and finals.
\item \textbf{The Digital Product Passport.} Registry infrastructure from July
2026; batteries mandatory 2027; steel, textiles, electronics phased
2027--2030\cite{espr}. DPP payloads are machine-readable product sustainability
data --- structurally the same output class CarbonPass already generates.
\item \textbf{Domestic escalation.} Threshold lowering, fee increases, ETS
pilot $\sim$2027\cite{ecobiz} --- each step converts more of Taiwan's
manufacturing SMEs from spectators into reporters, with CarbonPass's inventory
by-product as their on-ramp.
\item \textbf{The NDC 3.0 arithmetic.} A 38\%$\pm$2\% cut by 2035\cite{ndc}
cannot be delivered by $\sim$500 directly regulated entities alone; the state
needs measurement capillaries into the SME layer. A tool the government can
adopt through the Hackathon's own pipeline (Section~\ref{sec:fit}) is exactly
such a capillary.
\end{enumerate}

The expansion sequence is therefore: fasteners (2026, live compulsion) $\to$
downstream steel/aluminium categories (2028, pending) $\to$ DPP sectors
(2027--2030) $\to$ domestic carbon-fee entrants (2027 onward). Same engine,
widening catalogue.
